\section{Appendix E: Verification note}

The supplied theorem and lemma statements and their proofs are machine-checked in Lean 4 under their stated assumptions. The frozen formal layer covers identification and projection control in \cref{obj:lem:observable-odds,obj:lem:projection-control}; rational arithmetic and finite-budget realization in \cref{obj:lem:concrete-arithmetic-engine,obj:lem:fixed-budget-realization}; local implicit-function arguments, exact calibration, and model support in \cref{obj:lem:scalar-implicit-function,obj:lem:exact-calibrations,obj:lem:calibrated-support}; and original-record mixture bounds and the parametric floor in \cref{obj:lem:original-record-mixtures,obj:lem:parametric-null-floor}. The resulting coverage, expected-length comparison, and attaining construction are consolidated in \cref{obj:thm:finite-honest-upper,obj:thm:sharp-honest-frontier,obj:thm:full-honest-frontier-solution,obj:thm:full-honest-frontier-answer}.

The model restrictions defining \cref{obj:def:model} and the primitive contracts stated in the corresponding results are hypotheses of these implications. The global coverage criterion defining \cref{obj:def:honest-procedures} specifies the admissible statistical procedures. Admissibility of the concrete arithmetic engine, finite termination and endpoint enclosure for the represented construction, global coverage, and the sharp comparison for the length objective in \cref{obj:def:length-objective} are established conclusions. The recorded external-dependency sets for these declarations are empty.

For the represented construction, valid names satisfying \cref{obj:def:dyadic-name-convention} supply enclosing dyadic intervals for the public exponents and observed covariates. Together with the rational primitive contracts in \cref{obj:def:arithmetic-engine}, they support the rank selection and endpoint realization in \cref{obj:lem:fixed-budget-realization}. The explicit arithmetic witness in \cref{obj:lem:concrete-arithmetic-engine} supplies an admissible engine, and its combination with the dyadic-floor naming policy yields the concrete attaining interval identified in \cref{obj:thm:full-honest-frontier-solution}.

The formal source tree is
\texttt{CausalSmith/CausalSmith/Stat/STAT\_LogoddsLowsmoothFrontier\_Research/},
with its aggregate module at the corresponding \texttt{.lean} path. The
declaration-to-paper mapping is
\texttt{CausalSmith/doc/presentation/stat\_logodds\_lowsmooth\_frontier\_v1/presentation\_crosswalk.json};
it records repository source revision
\texttt{1efc342d4a68f9860ec0aa7692071d90a9230018}. In a checkout of that revision,
run
\begin{verbatim}
cd CausalSmith
lake build CausalSmith.Stat.STAT_LogoddsLowsmoothFrontier_Research
\end{verbatim}
The Lean toolchain is pinned in \texttt{CausalSmith/lean-toolchain}, and package
dependencies are pinned in \texttt{CausalSmith/lake-manifest.json} under the
project configuration \texttt{CausalSmith/lakefile.toml}. The same presentation
directory contains \texttt{lean\_snippets.json}, \texttt{references.bib},
\texttt{paper\_macros.tex}, and the crosswalk needed to reproduce the manuscript
links and inspect the exact declarations rendered in each numbered environment.
